% !TEX root = ../main.tex
\chapter*{Resumen}
\addcontentsline{toc}{chapter}{Resumen}

El sector retail se enfrenta al reto constante de equilibrar la disponibilidad de producto con la minimización de mermas. Este Trabajo de Fin de Grado presenta el desarrollo de un Sistema de Apoyo a la Decisión (DSS) End-to-End cuyo objetivo operativo es estimar y automatizar cuántas unidades debe pedir diariamente cada SKU bajo restricciones logísticas reales.

La metodología propuesta evoluciona el \textit{forecasting} clásico integrando modelos de Gradient Boosting con \textit{Mondrian Conformal Prediction} para garantizar la equidad estadística de la incertidumbre por categoría de producto. El sistema no se limita a predecir, sino que transforma estos pronósticos probabilísticos en decisiones de inventario mediante la formulación del \textit{Newsvendor}, que traduce la distribución de demanda estimada en una cantidad de pedido según la asimetría entre el coste de rotura y el de sobrestock. Finalmente, el modelo estadístico se ha operativizado bajo estándares MLOps, incluyendo explicabilidad automática (SHAP), trazabilidad de experimentos por artefactos versionados y un cuadro de mandos interactivo renderizado en servidor.

Los resultados empíricos demuestran que la corrección de la demanda censurada, bajo una evaluación justa que enfrenta a todas las estrategias contra una misma demanda real de referencia, reduce el coste logístico de inventario en un \pendiente{92{,}4}\% (con un incremento en el fill rate del \pendiente{60{,}2}\% al \pendiente{97{,}6}\%), y que, a coste de inventario prácticamente idéntico, el modelo probabilístico CatBoost eleva el nivel de servicio del \textbf{79,3\%} al \textbf{94,4\%} frente al modelo estacional ingenuo, pese a que este último obtiene el mejor MAE. Este enfoque integral contribuye directamente al \textbf{ODS 12 (Producción y Consumo Responsables)} al mitigar el desperdicio de productos perecederos mediante una cadena de suministro altamente eficiente, matemática e inteligible.

\textbf{Palabras clave:} Retail, Machine Learning, Investigación de Operaciones, Mondrian Conformal Prediction, Newsvendor, MLOps, ODS 12.

\newpage

\begin{otherlanguage}{catalan}
\chapter*{Resum}
\addcontentsline{toc}{chapter}{Resum}

El sector del retail s'enfronta al repte constant d'equilibrar la disponibilitat de producte amb la minimització de minves. Aquest Treball de Fi de Grau presenta el desenvolupament d'un Sistema de Suport a la Decisió (DSS) End-to-End amb l'objectiu operatiu d'estimar i automatitzar quantes unitats s'han de demanar diàriament per a cada SKU sota restriccions logístiques reals.

La metodologia proposada evoluciona el \textit{forecasting} clàssic integrant models de Gradient Boosting amb \textit{Mondrian Conformal Prediction} per a garantir l'equitat estadística de la incertesa per categoria de producte. El sistema no es limita a predir, sinó que transforma aquests pronòstics probabilístics en decisions d'inventari mitjançant la formulació del \textit{Newsvendor}, que tradueix la distribució de demanda estimada en una quantitat de comanda segons l'asimetria entre el cost de ruptura i el de sobreestoc. Finalment, el model estadístic s'ha operativitzat sota estàndards MLOps, incloent-hi explicabilitat automàtica (SHAP), traçabilitat d'experiments per artefactes versionats i un quadre de comandament interactiu renderitzat al servidor.

Els resultats empírics demostren que la correcció de la demanda censurada, sota una avaluació justa que enfronta totes les estratègies contra una mateixa demanda real de referència, redueix el cost logístic d'inventari en un \pendiente{92{,}4}\% (amb un increment del fill rate del \pendiente{60{,}2}\% al \pendiente{97{,}6}\%), i que, a un cost d'inventari pràcticament idèntic, el model probabilístic CatBoost eleva el nivell de servei del \textbf{79,3\%} al \textbf{94,4\%} respecte del model estacional ingenu, tot i que aquest darrer obté el millor MAE. Aquest enfocament integral contribueix directament a l'\textbf{ODS 12 (Producció i Consum Responsables)} en mitigar el malbaratament de productes peribles mitjançant una cadena de subministrament altament eficient, matemàtica i intel·ligible.

\textbf{Paraules clau:} Retail, Machine Learning, Investigació d'Operacions, Mondrian Conformal Prediction, Newsvendor, MLOps, ODS 12.
\end{otherlanguage}

\newpage

\begin{otherlanguage}{english}
\chapter*{Abstract}
\addcontentsline{toc}{chapter}{Abstract}

The retail sector faces the constant challenge of balancing product availability with waste minimization. This Bachelor's Thesis presents the development of an End-to-End Decision Support System (DSS) whose operational goal is to estimate and automate how many units of each SKU should be ordered daily under real logistical constraints.

The proposed methodology evolves classical forecasting by integrating Gradient Boosting models with Mondrian Conformal Prediction to ensure statistical fairness of uncertainty across product categories. The system goes beyond prediction by transforming probabilistic forecasts into inventory decisions via the \textit{Newsvendor} formulation, which turns the estimated demand distribution into an order quantity according to the asymmetry between stockout and overstock costs. Finally, the statistical model has been operationalized under MLOps standards, including automated explainability (SHAP), experiment traceability through versioned run artifacts, and an interactive server-rendered dashboard.

Empirical results demonstrate that correcting censored demand, under a fair evaluation that scores every strategy against one and the same real reference demand, reduces inventory logistics costs by \pendiente{92.4}\% (improving the fill rate from \pendiente{60.2}\% to \pendiente{97.6}\%), and that, at practically identical inventory cost, the probabilistic CatBoost model raises the service level from \textbf{79.3\%} to \textbf{94.4\%} against the seasonal naïve baseline, even though the latter attains the best MAE. This comprehensive approach directly contributes to \textbf{SDG 12 (Responsible Consumption and Production)} by mitigating perishable product waste through a highly efficient, mathematical, and intelligible supply chain.

\textbf{Keywords:} Retail, Machine Learning, Operations Research, Mondrian Conformal Prediction, Newsvendor, MLOps, SDG 12.
\end{otherlanguage}
